% =====================================================================
% MODELO PADRAO - TCC NO FORMATO SBC (Sociedade Brasileira de Computacao)
% =====================================================================
% Este e um modelo EM BRANCO, autocontido (nao depende de baixar
% sbc-template.sty a parte), reproduzindo a formatacao do template
% oficial da SBC: pagina A4, coluna UNICA do inicio ao fim (titulo,
% resumo, abstract e corpo do texto), sem divisao em colunas.
%
% Validado DIRETAMENTE contra o zip oficial baixado do site da SBC em
% www.sbc.org.br/documentosinstitucionais > Publicações > "Templates
% para Artigos e Capítulos de Livros" (modelosparapublicaodeartigos.zip,
% pasta Template_SBC/template-latex/sbc-template.tex + .sty). Reproduz
% aqui, sem depender do .sty externo: A4, coluna UNICA do inicio ao fim,
% margens 3,5/2,5/3/3 cm, fonte Times 12pt, SEM numeracao de pagina,
% titulo 16pt, autores/endereco 12pt, e-mail footnotesize monoespacado,
% legendas em fonte sem serifa 10pt com separador "." (nao ":"), e
% resumo/abstract em itálico com recuo simétrico de 0,8cm — tudo
% conforme o texto e o codigo-fonte do .sty oficial.
%
% COMO USAR:
%   1. Substitua todo texto entre colchetes [ ] pelo conteudo real.
%   2. Compile com pdflatex (rode duas vezes para os numeros de
%      secao/figura/tabela ficarem corretos):
%        pdflatex modelo-tcc-sbc.tex
%        pdflatex modelo-tcc-sbc.tex
%   3. Leia o "Manual de normas SBC para o TCC" antes de preencher —
%      ele explica quantidade de paginas, estrutura obrigatoria e
%      regras de citacao que este arquivo, sozinho, nao impõe.
% =====================================================================

\documentclass[a4paper,12pt]{article}

% ---------------------------------------------------------------------
% PACOTES
% ---------------------------------------------------------------------
\usepackage[utf8]{inputenc}
\usepackage[T1]{fontenc}
\usepackage[brazilian]{babel}
\usepackage{times}
\usepackage[a4paper,top=3.5cm,left=3cm,right=3cm,bottom=2.5cm]{geometry}
\usepackage{graphicx}
\usepackage{url}
\usepackage{indentfirst}
\usepackage{titlesec}
\usepackage[bf,sf,footnotesize,indent]{caption}

% ---------------------------------------------------------------------
% ESPACAMENTO DE PARAGRAFO (igual ao padrao SBC)
% ---------------------------------------------------------------------
\setlength{\parindent}{1.27cm}
\setlength{\parskip}{6pt}
\flushbottom

% ---------------------------------------------------------------------
% FORMATO DE SECOES: 13pt negrito (secao), 12pt negrito (subsecao),
% numero seguido de PONTO (ex.: "1. Introdução"), como no .sty oficial.
% ---------------------------------------------------------------------
\titleformat{\section}{\fontsize{13}{15}\bfseries}{\thesection.}{1em}{}
\titleformat{\subsection}{\fontsize{12}{14}\bfseries}{\thesubsection.}{1em}{}
\titleformat{\subsubsection}{\fontsize{12}{14}\bfseries\itshape}{\thesubsubsection.}{1em}{}
\titlespacing*{\section}{0pt}{12pt}{6pt}
\titlespacing*{\subsection}{0pt}{10pt}{4pt}

% Legendas: tabela acima, figura abaixo; fonte sem serifa, 10pt (footnotesize),
% negrito, separador "." (nao ":"), recuo de 0,8cm quando ocupar mais de 1 linha
% — exatamente como definido em sbc-template.sty.
\renewcommand{\captionfont}{\sffamily\footnotesize\bfseries}
\setlength{\captionmargin}{0.8cm}
\captionsetup{labelsep=period}

% ---------------------------------------------------------------------
% AMBIENTES RESUMO / ABSTRACT — recuo simetrico de 0,8cm e corpo em
% itálico, reproduzindo \newenvironment{resumo}/\abstract do .sty oficial.
% ---------------------------------------------------------------------
\newenvironment{resumo}{%
  \list{}{\advance\topsep by6pt\relax
    \leftmargin=0.8cm
    \labelwidth=0pt
    \listparindent=0pt
    \itemindent\listparindent
    \rightmargin\leftmargin}\item[\hskip\labelsep\bfseries\itshape Resumo.]\itshape}%
  {\endlist}

\renewenvironment{abstract}{%
  \list{}{\advance\topsep by6pt\relax
    \leftmargin=0.8cm
    \labelwidth=0pt
    \listparindent=0pt
    \itemindent\listparindent
    \rightmargin\leftmargin}\item[\hskip\labelsep\bfseries\itshape Abstract.]\itshape}%
  {\endlist}

\begin{document}
\pagestyle{empty} % a SBC exige numeracao de pagina SUPRIMIDA

% =====================================================================
% BLOCO DE TITULO, AUTORES, RESUMO E ABSTRACT
% =====================================================================
\begin{center}
  % Titulo em 16pt negrito (\fontsize{16}{20}), como no .sty oficial. O
  % modelo da SBC, em si, exige apenas o titulo no idioma do artigo; a
  % segunda linha (inglês) é um acréscimo institucional da UNIT, feito
  % em corpo menor para manter a hierarquia visual com o titulo oficial.
  {\fontsize{16}{20}\selectfont\bfseries [TÍTULO DO TRABALHO EM PORTUGUÊS]}\\[4pt]
  {\fontsize{14}{18}\selectfont\bfseries [Title of the Work in English]}\\[12pt]

  % Inclua o(a) Coorientador(a) apenas se houver (ver Manual de Normas —
  % Capítulo II, "Do Coorientador"). Nesse caso, acrescente-o logo após
  % o(a) Orientador(a); se for de outra instituição, use um novo número
  % de afiliação ($^{2}$) e replique o bloco de afiliação abaixo.
  {\normalsize\bfseries [Nome do(a) Autor(a)]$^{1}$, [Nome do(a) Orientador(a)]$^{1}$[, Nome do(a) Coorientador(a)]$^{1}$}\\[6pt]

  {\normalsize
    $^{1}$[Curso/Departamento — Nome da Instituição]\\{}
    [Cidade, Estado] -- [Sigla-UF]\\{}
    \footnotesize\texttt{\{email.autor,email.orientador\}@dominio.edu.br}
  }
\end{center}
\vspace{12pt}

\begin{resumo}
[Escreva aqui o resumo em português, em um único parágrafo, sem
citações bibliográficas, cobrindo objetivo, método, principais
resultados e conclusão do trabalho.]
\end{resumo}
\noindent\textbf{Palavras-chave:} [palavra 1], [palavra 2], [palavra 3].

\vspace{6pt}

\begin{abstract}
[Write here the English version of the abstract, following the same
structure as the Resumo above.]
\end{abstract}
\noindent\textbf{Keywords:} [keyword 1], [keyword 2], [keyword 3].

\vspace{18pt}

% =====================================================================
% CORPO DO TEXTO — coluna única, do mesmo jeito que o resto do artigo.
% =====================================================================

\section{Introdução}

[Contextualize o problema, apresente a motivação, o objetivo geral e
os objetivos específicos do trabalho, e feche o parágrafo final
descrevendo como o restante do texto está organizado
(ex.: "A Seção 2 apresenta os trabalhos relacionados... ").]

\section{Trabalhos Relacionados}

[Revisão da literatura / referencial teórico. Cite os trabalhos no
formato \cite{exemplo2024} e compare-os com a proposta deste TCC.]

\section{Metodologia}

[Descreva os materiais, métodos, ferramentas e procedimentos
utilizados para desenvolver o trabalho.]

\section{Desenvolvimento / Resultados}

[Apresente o que foi construído, experimentado ou analisado, e
discuta os resultados obtidos. Utilize figuras e tabelas conforme
necessário -- veja os exemplos comentados abaixo.]

% --- Exemplo de figura (remova o comentário e ajuste o caminho) ---
% \begin{figure}[ht]
%   \centering
%   \includegraphics[width=.7\linewidth]{caminho-da-imagem}
%   \caption{Descrição da figura.}
%   \label{fig:exemplo}
% \end{figure}

% --- Exemplo de tabela ---
% \begin{table}[ht]
%   \centering
%   \caption{Descrição da tabela.}
%   \label{tab:exemplo}
%   \begin{tabular}{lcc}
%     \hline
%     Coluna 1 & Coluna 2 & Coluna 3 \\ \hline
%     valor & valor & valor \\ \hline
%   \end{tabular}
% \end{table}

\section{Conclusão}

[Retome o objetivo do trabalho, resuma as principais contribuições e
resultados, aponte limitações e sugestões de trabalhos futuros.]

% =====================================================================
% DECLARACAO DE USO DE IA (obrigatoria — Manual de Normas, Art. 12)
% =====================================================================
\section*{Declaração de Uso de Inteligência Artificial}

[Conforme a Portaria CNPq nº 2.664/2026 e o Art. 12º do Manual de Normas
de TCC, declare aqui o uso -- ou a ausência de uso -- de ferramentas de
Inteligência Artificial (IA) generativa neste trabalho. Se utilizou
alguma ferramenta, informe: (i) o nome da ferramenta e, se aplicável, a
versão ou modelo; (ii) a etapa do trabalho em que foi empregada (ex.:
revisão textual, tradução, geração ou revisão de código-fonte, análise
de dados); e (iii) a finalidade específica de uso em cada etapa. Caso
nenhuma ferramenta de IA generativa tenha sido utilizada, declare
expressamente essa ausência de uso -- esta seção não pode ser omitida.]

% =====================================================================
% REFERENCIAS
% =====================================================================
% Ajuste o estilo de citação conforme o que a coordenação do curso
% exigir (numérico, como abaixo, ou autor-data no padrão ABNT). Veja
% o manual de normas para orientação.

\begin{thebibliography}{9}

\bibitem{exemplo2024}
Sobrenome, N. dos.
\newblock Título do trabalho citado como exemplo.
\newblock \emph{Nome do Periódico ou Evento}, v. 1, n. 1, p. 1--10, 2024.

\end{thebibliography}

\end{document}
